\section{ImageNet-9 Experimental Details}
\label{app:imagenet9-setup}

This appendix gives the construction and secondary diagnostics behind Section~\ref{sec:imagenet9}. The main text keeps only the tests needed for the central argument. Here we specify the data splits, model zoo, paired background interventions, independent behavioral indicators, magnitude-controlled evaluation, spatial-attribution baselines, and cross-architecture comparison.

\subsection{Data variants and disjoint splits}

We use the official ImageNet-9 Backgrounds Challenge variants~\citep{xiao2021noise}: Original, Mixed-Same, Mixed-Rand, Mixed-Next, Only-FG, and the challenge backgrounds. The variants preserve foreground identity while changing how the background is constructed. We match variants by foreground identifier and retain 4,050 foregrounds for which the required variants are available.

A deterministic split keeps the roles of the data separate. We reserve 1,644 paired foregrounds for the broad model-level response scan and an 820-foreground deep pool for the sample-level benchmark; 768 foregrounds from this pool are used for the expensive baseline comparison. The remaining foregrounds are not used in that sample-level benchmark. No foreground appears in more than one split.

\subsection{Model zoo and common score space}

The 72-model zoo contains two complementary groups. The first consists of 24 off-the-shelf ImageNet-1k classifiers from torchvision and timm. These models were trained independently of our audit and span ResNet, ResNeXt, WideResNet, DenseNet, EfficientNet, RegNet, ConvNeXt, DeiT, ViT, BEiT, Swin, MaxViT, and CoAtNet architectures.

The second group contains 48 models trained directly on ImageNet-9. We fine-tune six backbones---ResNet-50, ConvNeXt-Tiny, EfficientNet-B3, RegNetY-8GF, Swin-T, and ViT-B/16---on four versions of the training data: Original, Mixed-Rand, Mixed-Next, and Only-FG, using two random seeds. This $6\times4\times2$ design broadens the range of learned background dependence while retaining the same ImageNet-9 prediction task.

We train the fine-tuned models for eight epochs with AdamW, cosine decay, one warm-up epoch, BF16 computation, random resized crops, and horizontal flips. ResNet-50, EfficientNet-B3, and RegNetY-8GF use batch size 256 and learning rate $3\times10^{-4}$. ConvNeXt-Tiny, Swin-T, and ViT-B/16 use batch size 128 and learning rate $10^{-4}$.

All models are evaluated in the same nine-class score space. For an ImageNet-1k model, we compute its 1,000-way softmax and sum the probabilities belonging to each official ImageNet-9 superclass. We do not sum logits or apply a second softmax. The ImageNet-9 fine-tuned models output nine logits directly. In both groups, the scalar analyzed by \decaf is the probability assigned to the true ImageNet-9 superclass.

\subsection{Paired background interventions and reveal paths}

For each foreground, we construct two paired background interventions. Same--Rand compares Mixed-Same with Mixed-Rand: the foreground is unchanged, while a same-class background is replaced by a random-class background. This pair is used for the evidence and fragility analyses. Same--Next compares Mixed-Same with Mixed-Next, where the replacement background comes from superclass $(y+1)\bmod 9$. This pair creates a directional background change for the contradiction analysis.

The primary reveal begins from a shared blurred midpoint. For endpoints $\mathbf x^{+}$ and $\mathbf x^{-}$,
\begin{equation}
    \mathbf x^{0}
    =
    \operatorname{Blur}\!\left(
        \frac{\mathbf x^{+}+\mathbf x^{-}}{2}
    \right),
\end{equation}
and each branch is linearly revealed from $\mathbf x^{0}$ to its endpoint. We evaluate nine stages,
$t\in\{0,0.125,\ldots,1\}$. The main endpoint threshold is $\varepsilon=0.02$ in true-class probability.

To test whether conclusions depend on the reveal path, we also use nested-patch reveal on an $8\times8$ grid. Patches are ordered by endpoint-difference energy, with a fixed random tie-break. Two independently tie-broken orders are evaluated. Both branches use the same patch order and the same neutral starting image. Appendix~\ref{app:imagenet9-results} reports the corresponding path-order and threshold checks.

\subsection{Behavioral indicators defined without \decaf}
\label{app:imagenet9-labels}

The sample-level benchmark defines three observable behavioral indicators directly from model predictions. None uses $E$, $C$, or $F$, and the indicators are allowed to overlap.

\paragraph{Evidence behavior.}
Set $Y_E=1$ when Mixed-Same is classified correctly and replacing its background with Mixed-Rand either changes the predicted class or lowers the true-class probability by at least $0.20$. This records a consequential change caused by replacing the background.

\paragraph{Contradiction behavior.}
Set $Y_C=1$ when changing only the background makes the model switch from the foreground class to the class associated with the new Mixed-Next background. This is an independently observed directional reversal in the model's prediction.

\paragraph{Fragility behavior.}
We first restrict attention to Same--Rand pairs whose fully revealed endpoints produce little difference, $|d|<0.02$. The endpoint pair alone cannot tell us whether such a case is otherwise sensitive, so we test it separately with additional background changes that are not part of the \decaf reveal. These changes use Gaussian blur, Gaussian noise, pixelation, color shift, and patch shuffle, each at two severities. Set $Y_F=1$ if any of them changes the predicted class or shifts the true-class probability by at least $0.20$.

These indicators are external behavioral checks, not definitions of the \decaf components. Their purpose is to ask whether the response roles obtained from the paired trajectory agree with behavior measured independently of that trajectory decomposition.

\subsection{Magnitude-controlled evaluation}

The deep benchmark uses 32 models and 768 foregrounds, giving 24,576 model--image units. We use two complementary ways to control ordinary response magnitude.

First, we divide the units into 20 quantile bins of $\Abs$. Within a bin, a behavioral AUROC is computed only when both positive and negative examples are present for that indicator. The main text reports both the valid-bin summary and the stricter common-support comparison in which all three indicators are represented.

Second, we construct one-to-one matched comparisons between cases with different behavioral indicator patterns but nearly identical ordinary response magnitudes. We require the relative difference in $\Abs$ to be at most $5\%$, yielding 8,289 matched comparisons.

Because $Y_E$, $Y_C$, and $Y_F$ may overlap, the matched evaluation does not force every case into one ground-truth class. For a case with at least one active indicator, we examine the largest score among the three response-role coordinates. The case receives full credit when the largest coordinate corresponds to an active behavioral indicator. If several coordinates tie for the maximum, credit is divided across the tied coordinates. The reported matched-pair accuracy averages this role-agreement score over the matched cases. Ordinary magnitude has only one scalar value and therefore cannot prefer evidence, contradiction, or fragility.

The implementation also records a separate pairwise-ranking diagnostic: when an independent behavioral indicator differs across a matched pair, it asks whether the corresponding response-role score changes in the same direction. This ranking diagnostic is distinct from the role-agreement accuracy reported in Section~6.2.

\subsection{Spatial-attribution baselines}

The deep benchmark includes six standard spatial-attribution methods: Input$\times$Gradient, Integrated Gradients with 16 steps, SmoothGrad with 16 noise samples, BlurIG with 12 blur levels, a $7\times7$ Occlusion grid, and RISE with 256 masks. For each method, we sum absolute attribution inside the foreground--background difference region. These scores describe where the prediction is sensitive; they are not used to define the three behavioral indicators above.

\subsection{Cross-architecture comparison and spatial complementarity}

\begin{figure}[!t]
    \centering
    \includegraphics[width=0.5\linewidth]{figures/imagenet9/fig5d_loao_decodability_synergy.pdf}
    \caption{\textbf{Cross-architecture transfer and spatial complementarity.}
    All learned results hold out one architecture family during fitting.
    Response features from \decaf nearly match the supervised response-statistics reference, while adding spatial attribution provides a further gain and approaches the full supervised reference.}
    \label{fig:app-in9-cross-family}
\end{figure}

For the learned comparison, every representation is evaluated with the same supervised protocol: train on all but one architecture family and evaluate on the held-out family. The \decaf representation uses $(M,E,C,F)$. A response-statistics reference uses $(M,\Abs,\mathrm{Net},\mathrm{SignFlip},\mathrm{active\ rate})$. The spatial representation uses the six attribution scores above. We also evaluate \decaf combined with the spatial scores and a full reference combining response statistics with spatial attribution.

Figure~\ref{fig:app-in9-cross-family} contains the panel moved from the main text. Using response information alone, \decaf reaches leave-one-architecture-family-out macro-AUROC $0.862$, close to the supervised response-statistics reference at $0.867$. Spatial attribution alone reaches $0.737$. Combining \decaf with spatial attribution reaches $0.918$, close to the full supervised reference at $0.921$. This is best read as a complementarity check: the response decomposition and spatial attribution retain different information.



\subsection{Statistical reporting}

For direct behavioral comparisons, we report AUROC and AUPRC. For learned comparisons, we report leave-one-architecture-family-out results so that the evaluated architecture family is not used to fit the probe. Undefined AUROCs remain missing rather than being replaced by chance. Where bootstrap intervals are reported, models---rather than individual model--image rows---are the resampling unit.
